% answers_short.tex —— 2014 年·安徽卷·理科 · 简短答案速查
% 由 tools/make_short_answers.py 自 answers.tex 的 \gksolhead/\gkanswer 提取，
% 并与 2014/2014_数学真题_安徽理_答案要点.md 逐题比对后生成，勿手改（改源 MD 后重跑）。

\begin{gknotice}
\textbf{使用说明}：本节只给最终答案便于快速核对，完整推导见解析版。答案为编写者独立推导，\textbf{非官方标准答案}。
\end{gknotice}

\gkpart{选择题}
\begin{center}\begin{tabular}{ccccccccccc}
\toprule
题号 & 1 & 2 & 3 & 4 & 5 & 6 & 7 & 8 & 9 & 10 \\
\midrule
答案 & C & B & B & D & D & A & A & C & D & A \\
\bottomrule
\end{tabular}\end{center}

\gkpart{填空题}
\begin{itemize}[leftmargin=2.6em]
  \item[\textbf{11.}] $\dfrac{3\pi}8$。平移后相位为 $\dfrac\pi4-2\varphi$，令其等于 $\dfrac\pi2+k\pi$，最小正值 $\varphi=\dfrac{3\pi}8$。
  \item[\textbf{12.}] $1$。设 $A=a_1+1$，由 $(A+2d+2)^2=A(A+4d+4)$ 得 $d=-1$，三数相等，公比为 $1$。
  \item[\textbf{13.}] $3$。$a_0=1$，$a_1=\dfrac na=3$，$a_2=\dfrac{n(n-1)}{2a^2}=4$，联立得 $a=3$。
  \item[\textbf{14.}] $x^2+\dfrac{3y^2}{2}=1$。取 $A(c,b^2)$，直线参数代入椭圆，由两根 $1$ 与 $-\dfrac13$ 得 $b^2=2c^2$，结合 $c^2=1-b^2$ 得 $b^2=\dfrac23$。
  \item[\textbf{15.}] ②④。$S=k|\boldsymbol a-\boldsymbol b|^2+|\boldsymbol b|^2+4\boldsymbol a\cdot\boldsymbol b$，$k=0,1,2$；②中 $S_{\min}=|\boldsymbol b|^2$，④由 $|\boldsymbol b|(|\boldsymbol b|-4|\boldsymbol a|)>0$ 得证。
\end{itemize}

\gkpart{解答题}
\begin{itemize}[leftmargin=2.6em]
  \item[\textbf{16.}] (1) $a=2\sqrt3$。由 $a=6\cos B$ 与余弦定理 $9=a^2+1-\dfrac{a^2}3$ 解得；(2) $\dfrac{4-\sqrt2}{6}$。$\cos A=-\dfrac13$，$\sin A=\dfrac{2\sqrt2}3$，再用和角公式
  \item[\textbf{17.}] (1) $\dfrac{56}{81}$。甲2局、3局、4局获胜的概率分别为 $\dfrac49$、$\dfrac4{27}$、$\dfrac8{81}$，相加；(2) $P(X=2)=\dfrac59$，$P(X=3)=\dfrac29$，$P(X=4)=\dfrac{10}{81}$，$P(X=5)=\dfrac8{81}$；$E(X)=\dfrac{224}{81}$。按首次出现连胜的序列计数
  \item[\textbf{18.}] (1) 记 $x_1=\dfrac{-1-\sqrt{4+3a}}3$，$x_2=\dfrac{-1+\sqrt{4+3a}}3$；减区间为 $(-\infty,x_1)$、$(x_2,+\infty)$，增区间为 $(x_1,x_2)$。由 $f'(x)=-(3x^2+2x-1-a)$ 的符号得；(2) 最大值点当 $0<a<4$ 时为 $x=\dfrac{\sqrt{4+3a}-1}{3}$，当 $a\geqslant4$ 时为 $x=1$；最小值点当 $0<a<1$ 时为 $x=1$，当 $a=1$ 时为 $x=0$ 或 $x=1$，当 $a>1$ 时为 $x=0$。比较 $f(0)=1$、$f(1)=a$ 与极大值点
  \item[\textbf{19.}] (1) 见解析版；(2) $\dfrac{p_1^2}{p_2^2}$。三顶点坐标都随 $p$ 同比缩放，面积比为相似比的平方
  \item[\textbf{20.}] (1) 见解析版；(2) 上、下两部分的体积之比为 $11:7$。取标准梯形积分，下半体积为 $\dfrac{7H}{12}$，上半为 $\dfrac{11H}{12}$；(3) $\dfrac\pi4$。$S_{\triangle ACD}=4=\dfrac12\cdot CD\cdot d(A,CD)$ 得 $d(A,CD)=4$，故 $\tan\theta=\dfrac{A_1A}{d}=1$
  \item[\textbf{21.}] (1) 见解析版；(2) 见解析版
\end{itemize}
